%% usercommands.tex -- G. Metzler's user commands (from the template of the poly
%% "Statistical Supervised Machine Learning"), cleaned for the PAC-Bayes notes:
%% the duplicated block of the original file has been merged, and the packages
%% that are already loaded by polystyle.sty (or that clash with algorithm2e,
%% such as algorithmic) are not loaded again.

\newcommand{\todo}[1]{\textcolor{red}{\textbf{#1}}}
\definecolor{bleu}{RGB}{11,36,251}
\definecolor{vert}{RGB}{15,126,18}

% Packages
\usepackage{xifthen}
\usepackage{bbm}
\usepackage[makeroom]{cancel}
\usepackage{euscript}
\usepackage{mathrsfs}
\usepackage{amsfonts}
\usepackage{pgffor}
\usepackage{pgfplots}
\usepackage{pgfplotstable}
\pgfplotsset{compat=1.18}
\usepgfplotslibrary{fillbetween}
\usepackage{verbatim}
\usepackage{stmaryrd}
\usetikzlibrary{fit,plotmarks,decorations,arrows,shapes,backgrounds,matrix,chains,positioning,decorations.pathreplacing}

% Allowing page breaks in align
\allowdisplaybreaks

% Referencing a footnote
\makeatletter
\newcommand\footnoteref[1]{\protected@xdef\@thefnmark{\ref{#1}}\@footnotemark}
\makeatother

% Tables
\newcolumntype{L}[1]{>{\raggedright\let\newline\\\arraybackslash\hspace{0pt}}m{#1}}
\newcolumntype{C}[1]{>{\centering\let\newline\\\arraybackslash\hspace{0pt}}m{#1}}
\newcolumntype{R}[1]{>{\raggedleft\let\newline\\\arraybackslash\hspace{0pt}}m{#1}}

% Other commands
\newcommand{\dprime}{{\prime \prime}}
\newcommand{\tprime}{{\prime \prime \prime}}
\newcommand{\tR}[1]{\tRhelp#1\relax}
\def\tRhelp#1#2\relax{L_{\csname dom#1\endcsname#2}} % True Risk
\newcommand{\eR}[1]{\eRhelp#1\relax}
\def\eRhelp#1#2\relax{\hat{L}_{\csname set#1\endcsname#2}}
\newcommand{\loss}[1]{l\ifthenelse{\isempty{#1}{}}{}{\left(#1\right)}}
\newcommand{\algA}{\mathcal{A}}
\newcommand{\negspace}[1]{\phantom{\hspace*{-#1}}}

% Math operators
\DeclareMathOperator*{\argmin}{\arg\,\min}
\DeclareMathOperator*{\argmax}{\arg\,\max}
\DeclareMathOperator*{\expect}{\mathbb{E}}
\DeclareMathOperator*{\prob}{\mathbb{P}}
\DeclareMathOperator*{\var}{Var}

% Vectors
\newcommand{\mydefv}[1]{\expandafter\newcommand\csname v#1\endcsname{\mathbf{#1}}}
\newcommand{\mydefallv}[1]{\ifx#1\mydefallv\else\mydefv{#1}\expandafter\mydefallv\fi}
\mydefallv abdcefilmopquvwxyz\mydefallv % \va, \vb, ..., \vz

% Vectors symbols
\newcommand{\mydefvsym}[1]{\expandafter\newcommand\csname v#1\endcsname{\boldsymbol{\csname #1\endcsname}}}
\newcommand{\mydefallvsym}[1]{\ifx#1\mydefallvsym\else\mydefvsym{#1}\expandafter\mydefallvsym\fi}
\mydefallvsym {alpha}{beta}{delta}{gamma}{sigma}{theta}{varepsilon}{xi}\mydefallvsym % \valpha, ...
\newcommand{\bs}[1]{\bm{#1}}

% Matrices
\newcommand{\mydefm}[1]{\expandafter\newcommand\csname m#1\endcsname{\mathbf{#1}}}
\newcommand{\mydefallm}[1]{\ifx#1\mydefallm\else\mydefm{#1}\expandafter\mydefallm\fi}
\mydefallm ABCDEFGHIJKLMNOPQRSTUVWXYZ\mydefallm % \mA, ..., \mZ

% Matrices symbols
\newcommand{\mydefmsym}[1]{\expandafter\newcommand\csname m#1\endcsname{\boldsymbol{\csname #1\endcsname}}}
\newcommand{\mydefallmsym}[1]{\ifx#1\mydefallmsym\else\mydefmsym{#1}\expandafter\mydefallmsym\fi}
\mydefallmsym {Sigma}{Gamma}{Phi}\mydefallmsym

% Domain
\newcommand{\mydefdom}[1]{\expandafter\newcommand\csname dom#1\endcsname{\mathscr{#1}}}
\newcommand{\mydefalldom}[1]{\ifx#1\mydefalldom\else\mydefdom{#1}\expandafter\mydefalldom\fi}
\mydefalldom ABCDEFGLMN\mydefalldom

% Set
\newcommand{\mydefset}[1]{\expandafter\newcommand\csname set#1\endcsname{{#1}}}
\newcommand{\mydefallset}[1]{\ifx#1\mydefallset\else\mydefset{#1}\expandafter\mydefallset\fi}
\mydefallset BNPSTV\mydefallset

% Distribution over a domain
\newcommand{\mydefdistr}[1]{\expandafter\newcommand\csname distr#1\endcsname{\mathcal{D}_{\csname dom#1\endcsname}}}
\newcommand{\mydefalldistr}[1]{\ifx#1\mydefalldistr\else\mydefdistr{#1}\expandafter\mydefalldistr\fi}
\mydefalldistr CDNSTV\mydefalldistr

% Space
\newcommand{\mydefspace}[1]{\expandafter\newcommand\csname space#1\endcsname{\mathcal{#1}}}
\newcommand{\mydefallspace}[1]{\ifx#1\mydefallspace\else\mydefspace{#1}\expandafter\mydefallspace\fi}
\mydefallspace BCDEFGHIKLMNOPRSUVXYZ\mydefallspace

% Function symbols
\newcommand{\mydeffsym}[1]{\expandafter\newcommand\csname f#1\endcsname[2][]{\csname #1\endcsname##1\ifthenelse{\equal{##2}{}}{}{\!\left(##2\right)}}}
\newcommand{\mydefallfsym}[1]{\ifx#1\mydefallfsym\else\mydeffsym{#1}\expandafter\mydefallfsym\fi}
\mydefallfsym {phi}{epsilon}{eta}\mydefallfsym

% Number Set
\newcommand{\mydefnset}[1]{\expandafter\newcommand\csname nset#1\endcsname{\mathbb{#1}}}
\newcommand{\mydefallnset}[1]{\ifx#1\mydefallnset\else\mydefnset{#1}\expandafter\mydefallnset\fi}
\mydefallnset CKNPQRSZ\mydefallnset

% Norm
\newcommand{\normTwo}[1]{\left\|#1\right\|_2}
\newcommand{\normF}[1]{\left\|#1\right\|_\mathcal{F}}
\newcommand{\normOne}[1]{\left\|#1\right\|_1}
\newcommand{\normNuc}[1]{\left\|#1\right\|_\mathcal{*}}
\newcommand{\normMax}[1]{\left\|#1\right\|_\infty}
\newcommand{\norm}[1]{\left\|#1\right\|}
\newcommand{\abs}[1]{\left|#1\right|}
\newcommand{\hinge}[1]{\left[#1\right]_+}
\newcommand{\dotprod}[2]{\left\langle#1,#2\right\rangle}
\newcommand{\trace}[1]{\text{Tr}\left(#1\right)}
\newcommand{\inv}[1]{\left( #1 \right)^{-1}}

% Ceil and Floor
\newcommand{\ceil}[1]{{\left\lceil #1 \right\rceil}}
\newcommand{\floor}[1]{{\left\lfloor #1 \right\rfloor}}

% Useful to explain huge equations: a grey arrow and a comment between two
% lines of an align environment, as in the slides of the course
\newcommand{\cause}[1]{&{\color{lightgray}\downarrow{}\;\text{\small #1}} \nonumber\\}
\newcommand{\causee}[1]{&{\color{lightgray}\phantom{\downarrow{}}\;\text{\small #1}} \nonumber\\}
\newcommand{\underb}[1]{\lefteqn{{\color{lightgray}\underbrace{\phantom{#1}}_{=0}}}#1}
\newcommand{\NOcause}[1]{}

% Coloured under-braces and terms (blue, red, green, brown)
\newcommand{\umark}[1]{\lefteqn{{\color{bfig}\underbrace{\phantom{#1}}}}#1}
\newcommand{\omark}[1]{{\color{bfig}{#1}}}
\newcommand{\umarkk}[1]{\lefteqn{{\color{rfig}\underbrace{\phantom{#1}}}}#1}
\newcommand{\omarkk}[1]{{\color{rfig}{#1}}}
\newcommand{\umarkkk}[1]{\lefteqn{{\color{ForestGreen}\underbrace{\phantom{#1}}}}#1}
\newcommand{\omarkkk}[1]{{\color{ForestGreen}{#1}}}
\newcommand{\umarkkkk}[1]{\lefteqn{{\color{brown}\underbrace{\phantom{#1}}}}#1}
\newcommand{\omarkkkk}[1]{{\color{brown}{#1}}}
\newcommand{\inn}[2]{\langle#1,#2\rangle}

%% Additions for the PAC-Bayes notes: labelled coloured braces, i.e.
%% \ubrace{colour}{expression}{label} puts a coloured under-brace with a label.
\newcommand{\ubrace}[3]{{\color{#1}\underbrace{{\color{black}#2}}_{\text{\scriptsize\color{#1}#3}}}}
\newcommand{\obrace}[3]{{\color{#1}\overbrace{{\color{black}#2}}^{\text{\scriptsize\color{#1}#3}}}}
